% META: {"id":"1SPE-SUITES-EX-028","chapitre":"1SPE-SUITES","type_objet":"exercice","capacites":["1SPE-SUITES-C2","1SPE-SUITES-C5"],"capacites_codes":["C2","C5"],"methodes":["M2","M5"],"parcours":2,"competences":["calculer","raisonner"],"duree_min":20,"mode_creation":"ex_nihilo","sources_inspiration":[],"parametres_sympy":{"a":{"domaine":"entier","min":1,"max":5},"b":{"domaine":"entier","min":1,"max":20}},"fichier_tex":"chapitres/1SPE-SUITES/exercices/1SPE-SUITES-EX-028.tex","corrige_tex":"chapitres/1SPE-SUITES/corriges/1SPE-SUITES-CO-028.tex","status":"approved","origin":"generated","status_history":["generated","audited","accepted","approved"],"release_acceptance":"RELEASE_OWNER_BATCH_ACCEPTANCE","acceptance_closure_digest":"sha256:9b3ccf9a81c5520fbb7e03b7b2d3e7bf2057a02834b6d908bf3be5f49f3b553f","acceptance_receipt_digest":"sha256:4fad86f0282e0fa4e0419cca1ad6379ae7af5a280c04a4a90880f90a1c044242"}
% BEGIN-VERIFY
% from sympy import symbols, simplify
% n = symbols('n', integer=True, nonnegative=True)
% # u_n = -4n + 30, r = -4, u0 = 30
% u = -4*n + 30
% assert u.subs(n, 0) == 30
% assert u.subs(n, 1) == 26
% diff = u.subs(n, n+1) - u
% assert simplify(diff) == -4
% # Suite decroissante car r = -4 < 0
% # u_n > 0 => -4n + 30 > 0 => n < 7.5 => n <= 7
% assert u.subs(n, 7) == 2
% assert u.subs(n, 8) == -2
% END-VERIFY
\begin{exercice}{1SPE-SUITES-EX-028}{2}{20}
Soit $(u_n)$ la suite définie pour tout entier naturel $n$ par :
\[
  u_n = -4n + 30.
\]

\begin{enumerate}
  \item Calculer $u_0$, $u_1$ et $u_2$.

  \item Démontrer que $(u_n)$ est une suite arithmétique. Préciser sa raison et son premier terme.

  \item Étudier le sens de variation de $(u_n)$ en utilisant la raison. Justifier rigoureusement.

  \item Déterminer les valeurs de $n$ pour lesquelles $u_n > 0$, puis celles pour lesquelles $u_n \leq 0$.

  \item Soit la suite $(v_n)$ définie par $v_n = u_n^2$. Les termes $v_0$, $v_1$, $v_2$ sont-ils dans le même ordre croissant ou décroissant ? Expliquer pourquoi $(v_n)$ n'est pas monotone.
\end{enumerate}
\end{exercice}
